% ===========================================================================
%  Chapter 3 -- Results
%  Figures for this chapter go in figures/chapter3/
% ===========================================================================
\chapter{Results}\label{ch:results}
\chapterfigs{chapter3}

\section{Some plots}\label{sec:sfms}

\Cref{fig:sfms} uses the \texttt{subfigure} environment (package \texttt{subcaption}): each panel has its own \verb|\caption| and \verb|\label|, and the main caption describes the whole figure. Panels can be referenced individually: \cref{fig:sfms-cluster} shows a lower normalisation, as reported previously \citep[e.g.,][]{Brinchmann2004}.

\begin{figure}[tbp]
  \centering
  \begin{subfigure}[t]{0.48\textwidth}
    \centering
    \includegraphics[width=\textwidth]{sfms_field_mock}
    \caption{Field galaxies.}
    \label{fig:sfms-field}
  \end{subfigure}\hfill
  \begin{subfigure}[t]{0.48\textwidth}
    \centering
    \includegraphics[width=\textwidth]{sfms_cluster_mock}
    \caption{Cluster galaxies.}
    \label{fig:sfms-cluster}
  \end{subfigure}
  \caption{Star-forming main sequence (mock data) for (\subref{fig:sfms-field}) field and (\subref{fig:sfms-cluster}) cluster galaxies. Solid lines are linear fits.}
  \label{fig:sfms}
\end{figure}

